\subsection{Semantički SVM kao podsustav klasifikacije}

\paragraph{Motivacija.}
Uvođenje ovoga podsustava polazi od pretpostavke da se znatan dio AI-generiranih komentara ne razlikuje samo po stilu pisanja, nego i po semantičkoj namjeni. Takvi su komentari često oblikovani tako da pojačaju angažman: izazovu ljutnju, prodube političku ili ideološku podjelu, zaoštre raspravu ili potaknu niz daljnjih reakcija. Zbog toga je razumno očekivati da se u prostoru značenja pojavljuju obrasci koji nisu slučajni i koje je moguće geometrijski razdvajati. Ne tvrdi se da svaki AI komentar nužno ima takvu funkciju, nego da je taj signal dovoljno čest da ga vrijedi zasebno modelirati. U tom smislu ovaj podsustav ne traži samo pojedine riječi, nego pokušava razdvojiti komentare prema njihovoj semantičkoj ulozi.

Semantički SVM pridružuje komentaru $t$ vrijednost

\[
p_{\mathrm{sem}}(t) \approx P(y=1\mid t),
\]

pri čemu $y=1$ označava razred \texttt{AI}, a $y=0$ razred \texttt{HUMAN}. Ta se vrijednost dalje koristi kao jedan od triju ulaza završnog modela.

\paragraph{Vektorski prikaz i predobrada.}
Tekst se najprije prevodi u semantički vektorski prikaz $x$, zatim normira

\[
\hat{x}=\frac{x}{\|x\|_2},
\]

a potom standardizira po koordinatama

\[
z_j=\frac{\hat{x}_j-\mu_j}{\sigma_j}.
\]

L2-normiranje uklanja razlike u duljini vektora, dok standardizacija usklađuje mjerila pojedinih koordinata prije učenja linearnog razdjelnika.

Takav je slijed važan jer SVM ne radi nad rečenicama kao nizovima znakova, nego nad točkama u prostoru $\mathbb{R}^d$. Ugrađivački model već je prije toga proveo nelinearno semantičko preslikavanje teksta u vektor, pa klasifikator više ne razdvaja ``riječi'', nego geometrijske položaje komentara u dobivenom prostoru.

\paragraph{Zašto je linearan razdjelnik ovdje prirodan izbor.}
Nakon ugradnje tekst više nije prikazan u izvornom prostoru, nego u prostoru semantičkih značajki. Složeniji, nelinearni dio posla već je obavio model za vektorske prikaze. Zbog toga linearan SVM u tom novom prostoru često daje vrlo dobar kompromis između interpretabilnosti, brzine i točnosti. Drugim riječima, nelinearnost je sadržana u prikazu, a ne u samoj završnoj klasifikacijskoj funkciji.

\paragraph{Odlučna funkcija i geometrija margine.}
Nad vektorom $z$ uči se linearna odlučna funkcija

\[
f(z)=w^\top z+b.
\]

Granica razdvajanja zadana je jednadžbom $w^\top z+b=0$, a rubovi margine jednadžbama $w^\top z+b=\pm 1$. Budući da širina margine iznosi

\[
\frac{2}{\|w\|_2},
\]

minimizacija člana $\frac{1}{2}\|w\|_2^2$ ekvivalentna je maksimizaciji margine. Model, dakle, traži što širi pojas razdvajanja dvaju razreda.

To se može pokazati izravno. Neka su $z_+$ i $z_-$ točke na dvama rubovima margine, pa vrijedi

\[
w^\top z_+ + b = 1,
\qquad
w^\top z_- + b = -1.
\]

Oduzimanjem dobivamo

\[
w^\top (z_+ - z_-) = 2.
\]

Najkraća udaljenost dviju paralelnih hiperravnina mjeri se u smjeru njihova normale $w$, pa je $z_+ - z_- = \lambda w$ za neki skalar $\lambda$. Uvrštavanjem slijedi

\[
w^\top(\lambda w)=\lambda\|w\|_2^2=2,
\qquad
\lambda=\frac{2}{\|w\|_2^2}.
\]

Stoga je udaljenost između rubova margine

\[
\|z_+-z_-\|_2 = |\lambda|\,\|w\|_2 = \frac{2}{\|w\|_2}.
\]

Upravo zato minimizacija norme vektora $w$ znači širenje margine.

\paragraph{Optimizacija meke margine.}
Za oznake $\tilde y_i\in\{-1,+1\}$ rješava se problem

\[
\min_{w,b,\xi}\;\frac{1}{2}\|w\|_2^2 + C\sum_{i=1}^n \xi_i
\]

uz ograničenja

\[
\tilde y_i(w^\top z_i+b)\ge 1-\xi_i,
\qquad
\xi_i\ge 0.
\]

Polazište je tvrda margina, gdje bi se zahtijevalo $\tilde y_i(w^\top z_i+b)\ge 1$ za svaki primjer. Budući da stvarni tekstualni podaci nisu savršeno razdvojivi, uvodi se meka margina preko varijabli $\xi_i$.

Izraz $\tilde y_i(w^\top z_i+b)$ objedinjuje oba razreda u jednom uvjetu: pozitivan je kada je primjer na ispravnoj strani granice, a negativan kada je pogrešno razvrstan. Time jedno ograničenje istodobno pokriva i AI i HUMAN razred. Varijabla $\xi_i$ mjeri povredu idealne margine, pa vrijedi: $\xi_i=0$ za ispravno razvrstane primjere izvan margine, $0<\xi_i<1$ za primjere unutar margine, $\xi_i=1$ za primjere točno na granici odluke, a $\xi_i>1$ za pogrešno razvrstane primjere. Parametar $C$ određuje odnos između širine margine i kazne za pogreške: veći $C$ snažnije kažnjava povrede, a manji $C$ dopušta pravilniju, ali fleksibilniju granicu.

Isti se problem može zapisati i preko hinge-gubitka

\[
\ell_{\mathrm{hinge}}(\tilde y,f(z)) = \max\{0,1-\tilde y f(z)\},
\]

pa optimizacijska zadaća poprima oblik

\[
\min_{w,b}\;\frac{1}{2}\|w\|_2^2 + C\sum_{i=1}^n \max\{0,1-\tilde y_i(w^\top z_i+b)\}.
\]

Taj zapis jasno pokazuje zašto primjeri daleko od granice ne pridonose gubitku, dok primjeri unutar margine i pogrešno razvrstani primjeri nose kaznu.

\paragraph{Dualna formulacija i računalna provedba.}
U dualnoj formulaciji vrijedi

\[
w=\sum_{i=1}^n \alpha_i \tilde y_i z_i,
\]

pa na granicu razdvajanja bitno utječu samo primjeri s $\alpha_i>0$, odnosno potporni vektori. Oni su geometrijski oni primjeri koji ``nose'' razdvajanje: leže na rubu margine ili narušavaju njezine uvjete. Ta formulacija objašnjava zašto o svim primjerima ne ovisi završna granica jednako.

U klasičnom kernel-SVM pristupu dualni se problem često rješava metodama poput SMO postupka, jer je tada prirodno raditi s unutarnjim produktima među primjerima. U stvarnoj provedbi ovoga podsustava to nije slučaj: ovdje se koristi \texttt{LinearSVC} nad standardiziranim semantičkim prikazima. Zato je dualna slika važna za matematičko razumijevanje modela, ali nije doslovan opis numeričkog postupka koji se izvršava u kodu.

\paragraph{Vjerojatnosna kalibracija.}
SVM izravno daje potpisanu vrijednost odlučne funkcije

\[
s=f(z)=w^\top z+b,
\]

a ne vjerojatnost. Zbog toga se primjenjuje Plattova kalibracija

\[
P(y=1\mid s)\approx \frac{1}{1+\exp(As+B)}.
\]

To ima jasnu interpretaciju: velike pozitivne vrijednosti $s$ trebaju odgovarati vjerojatnostima bliskima $1$, velike negativne vrijednosti vjerojatnostima bliskima $0$, a prijelaz oko granice odvija se glatko, bez skokova. Parametri $A$ i $B$ određuju nagib i položaj te sigmoidne preslikbe.

Parametri $A$ i $B$ uče se isključivo na skupu za učenje, i to trostrukom unakrsnom kalibracijom: u svakom presjeku bazni \texttt{LinearSVC} trenira se na dvjema trećinama podataka, na preostaloj trećini računaju se vrijednosti odlučne funkcije, a na njima se zatim prilagođava sigmoidna funkcija. Time se izbjegava kalibriranje na istim vrijednostima na kojima je razdjelnik upravo naučen. Završna procjena dobiva se kao prosjek triju kalibriranih vjerojatnosti. Razvojni skup pritom ostaje odvojen i služi samo za vrednovanje.

\paragraph{Izlaz podsustava.}
Za pojedini komentar podsustav vraća kalibriranu vjerojatnost $p_{\mathrm{sem}}(t)$. Kad bi se koristio samostalno, binarna bi odluka mogla biti zadana pravilom

\[
\hat y = \mathbf{1}\{p_{\mathrm{sem}}(t)\ge 0.5\}.
\]

U ovom radu to nije središnje, jer završni model koristi samu vjerojatnost, a ne unaprijed zadani prag.

\paragraph{Rezultati na razvojnom skupu.}
Na punom skupu za učenje, uz unaprijed izračunate semantičke vektorske prikaze, dobivene su vrijednosti točnosti $0.90948$, preciznosti $0.90341$, odziva $0.91283$, mjere $F_1=0.90810$ i pokazatelja $\mathrm{ROC\ AUC}=0.96885$. Konfuzijska matrica na razvojnom skupu iznosi

\[
\begin{array}{c|cc}
 & \hat y=0 & \hat y=1 \\
\hline
y=0 & 19549 & 2022 \\
y=1 & 1806 & 18912
\end{array}
\]

što pokazuje da podsustav dobro razdvaja razrede i daje pouzdanu procjenu vjerojatnosti AI razreda.

\paragraph{Mogućnost jednostavnog poboljšanja.}
Najizravnije poboljšanje ovoga podsustava jest zamjena sadašnjeg modela za izradu semantičkih vektorskih prikaza snažnijim modelom, uz zadržavanje iste SVM i kalibracijske sheme. Drugim riječima, ne mijenja se klasifikator, nego kvaliteta prikaza na kojem on radi. U sadašnjoj izvedbi korišten je \texttt{all-MiniLM-L6-v2} zbog brzine, dok bi prirodan sljedeći korak bio \texttt{BAAI/bge-small-en-v1.5}, a uz više računalnih resursa i veći modeli iste namjene. Takva je izmjena tehnički jednostavna, ali računski skuplja, jer najveći dio vremena i dalje otpada na izradu semantičkih vektorskih prikaza.